% Appendix B. Soft-sensor reframe 2026-11-04 (branch docs/cep-soft-sensor-frame): moved from hardware.tex
% (sec:recording), detection.tex, sizing.tex and csd.tex of docs/cep-grounded-draft; text condensed, notation and
% terminology unified; figs 07, 09, 11, 14 cut; review round 2: near/far implementation and 1080p detail cut. Numbers: scripts/rerun_paper_numbers.py (docs/optics/rerun_results.md),
% scripts/detector_checks.py sweep. Claim changes: papers/cep/CLAIM_CHANGES.md.
\section{Classical detectors and near/far disambiguation}\label{app:detectors}

The pipeline that preceded this work used two classical detectors and a size prior to label each detection near or
far of the focal plane. This appendix records what the failed classical-detector baseline of \secref{sec:csd-time}
needs: how the detectors work, the depth they admit and what their output follows.

\subsection{Case-study clip}\label{sec:recording}

The frame-level analyses below use 60 frames (every 15th) of the 1080p screen recording \code{2026-02-05 10-36-32.mp4},
whose date is disputed and whose strike phase is unknown (\code{scripts/check_clip_overlap.py}), with sizes
area-equivalent at the modelled geometric-mean scale of 1.701~\um{} per clip pixel; they are detector outputs, not
crystal sizes.

\subsection{Edge and watershed detectors}\label{sec:edge}

A per-pixel sharpness map (log of the Gaussian-smoothed squared Laplacian, \code{annotator.sharpness}) and the optical
flow to the previous frame are combined into a score that is quantised into $n_\ell = 3$ layers; each detector runs
once per layer. The layers order pixels by sharpness and motion and are not depth measurements. Both detectors start
from the same Gaussian adaptive threshold (block 11~px, offset 2), which marks the dark crystal outlines on the bright
background. The \emph{edge detector} (\code{annotator.detection}) closes the outline map morphologically with an
elliptical kernel that grows by 4~px per layer from 3~px, fits external contours with minimum-area rotated rectangles,
and keeps a contour if its area lies in the size window and its aspect ratio is at most~3. The \emph{watershed
detector} (\code{optics.detect_watershed}) opens the outline map with a $3 \times 3$ kernel and distance-transforms its
complement. Seeds are the h-maxima ($h_\mathrm{max} = 2.5$~px) of the distance map that lie at least 6~px from any
outline and at least 12~px from each other (module constants of \code{optics.detect_watershed}, a choice). A marker-controlled watershed on the negative
distance map, a standard image-segmentation method~\citep{Vincent1991,BeucherMeyer1993}, grows one basin per seed;
each basin is fitted and filtered as above. A detection's size is the geometric mean of its rectangle sides.


\label{sec:window}The area window is read from \code{configs/camera_100.toml} by
\code{optics.config.detection_area_px}. Its lower bound is $3^2 = 9$ samples$^2$ (\code{min_side_px} $= 3$, choice);
its upper bound is the area of a 1~mm square at the focal plane, $\mathit{max\_size}^2 / (\mu_\mathrm{x} \mu_\mathrm{y})$, where the 1~mm
(\code{max_size_um}, choice) rounds up the 0.821~mm largest crystal of the first microscope reading. On the 1080p clips
the window is $[9,\ 345\,389]$~px$^2$, an area-equivalent 5--1000~\um; on the PAL recording it is
$[9,\ 68\,811]$~samples$^2$, where three samples are $12.0 \times 10.9$~\um. The earlier fixed window
$[80,\ 8000]$~px$^2$ corresponds on the clips to an area-equivalent 15--150~\um.

\subsection{Depth range admitted and the near/far question}\label{sec:selectivity}

Both detectors need dark outlines a few samples wide, so they admit only crystals near the focal plane. If a detector
loses an outline once the CoC exceeds five PAL samples (an assumption, not a measured threshold), the F/1.6 lens
admits $\delta = \pm 48$~\um{} around focus (Eq.~\ref{eq:coc}); counting in clip pixels gives $\pm 22$~\um.

\label{sec:ambiguity}Defocus blur is symmetric in reciprocal object distance (\ref{sec:coc}), so two depths
$d_\mathrm{near}$ and $d_\mathrm{far}$ with equal blur cannot be told apart from the blur. The apparent size of a crystal
of physical width $w_\mathrm{c}$, $w_\mathrm{c} / \mu(d_o)$, is however monotone in $d_o$, because the scale grows
with object distance (Eq.~\ref{eq:pixelsize}). For a near-side defocus $|\delta| \ll d_\mathrm{focus}$,
\begin{equation}\label{eq:ratio}
\frac{w_\mathrm{c}/\mu(d_\mathrm{near})}{w_\mathrm{c}/\mu(d_\mathrm{far})} - 1 = \frac{d_\mathrm{far}}{d_\mathrm{near}} - 1
\approx \frac{2|\delta|}{d_\mathrm{focus}}.
\end{equation}
Over the admitted $\pm 48$~\um{} the ratio departs from one by at most 0.17\,\% (0.08\,\% in clip pixels), whereas
5\,\% needs $|\delta| \approx 1.3$~mm. This bounds the magnification change used in
\secref{sec:layer}. It also answers the near/far question of the earlier pipeline: a size prior cannot separate
hypotheses that differ by 0.17\,\%, and on the clip 99.6--99.9\,\% of the off-focus detections had a near probability
within $0.500 \pm 0.002$, so its near/far labels were tie-breaks. The mean depth correction was 0.005--0.013\,\% of the
median size (\code{scripts/rerun_paper_numbers.py}).


\subsection{Detector output on the clip}\label{sec:csd-edge}

On 60 frames of the clip (\code{scripts/rerun_paper_numbers.py}), sizes are detector outputs, not crystal sizes.
Widening the window from 15--150~\um{} to 5--1000~\um{} raises the edge detections from 18\,823 to 47\,414, halves the
median from 0.037 to 0.019~mm and raises the CV from 0.83 to 2.16: the edge-detector distribution depends on the area
window. By visual reading, not checked against labels, the small detections are outline fragments and the large ones
rectangles around blurred regions.

\label{sec:watershed}The watershed output hardly changes with the window (52\,302 against 52\,784 detections, median
0.076~mm under both). Varying its seed constants directly (\tabref{tab:ws-sweep}; \code{scripts/detector_checks.py
sweep}, six single frames, so the configured median is 89 rather than 75~\um) shows that the kept basins cover
86\,\% of the cropped frame: the detector tiles the picture whether or not a crystal outline separates the basins. The
median hardly moves with the seed spacing between 6 and 16~px, but it follows the seed clearance, from 64 to 109~\um{}
for 3 to 9~px, and from 61 to 150~\um{} when both constants are scaled together. A second sweep on three frames with the
three sharpness layers pooled (\code{scripts/watershed_seed_sweep.py}) gives median basins of 36.5, 45.2 and
90.3 clip pixels for the constants 6/3, 12/6 and 24/12~px. A narrower CV than the edge
detector's is therefore not evidence of a narrower crystal distribution.


\begin{table}
\centering\tabfont
\caption{Watershed output on six frames of the case-study clip against its seed constants (configured window;
\texttt{scripts/detector\_checks.py sweep}; configured constants 12/6~px). Coverage: share of the cropped frame
inside kept basins.}\label{tab:ws-sweep}
\begin{tabular}{@{}rrrrr@{}}
\toprule
Spacing (px) & Clearance (px) & Per frame & Median (\textmu m) & Coverage \\
\midrule
6 & 6 & 693 & 88 & 86\,\% \\
12 & 6 & 680 & 89 & 86\,\% \\
24 & 6 & 565 & 100 & 85\,\% \\
12 & 3 & 1304 & 64 & 91\,\% \\
12 & 9 & 335 & 109 & 73\,\% \\
6 & 3 & 1363 & 61 & 91\,\% \\
24 & 12 & 166 & 150 & 59\,\% \\
\bottomrule
\end{tabular}
\end{table}

\subsection{First reading with an untuned Cellpose model}\label{sec:cellpose-first}

A first comparison of Cellpose outlines, not checked against labels, on frames of the 2026-10-30 recording with the
microscope photographs at matched times shows the camera sizes levelling off: the camera 90th percentile rises against
the microscope one with a log--log slope of only 0.32--0.51 (b1 0.32--0.42, b2 0.48--0.51 over four diameter settings), and the level it
settles at follows the detector's diameter setting (\code{scripts/cellpose_vs_lab.py}). Late in strike b1 the raw frames show crystals of about
90--190 PAL samples across, comparable at the modelled scale to the microscope crystals; flat-field division and
contrast-limited adaptive histogram equalisation split them into facets of 30--50 samples, which Cellpose outlines as
180--443 crystals per frame. Up to 11:46 the crystals span 3--9 samples and the
outlines follow texture. These failure modes motivate the fine-tuning of \secref{sec:detector}.
